\section*{Приложение А}
\addcontentsline{toc}{section}{Приложение А — Исходный код программы}

\begin{center}
	\Large\textbf{Исходный код программной реализации продукционной системы Маркова}
\end{center}

\lstset{
	language=Python,
	basicstyle=\ttfamily\footnotesize,
	keywordstyle=\color{blue},
	commentstyle=\color{codegreen},
	stringstyle=\color{red},
	showstringspaces=false,
	breaklines=true,
	frame=single,
	tabsize=4
}

\begin{lstlisting}
	import os
	import time
	import tkinter as tk
	from tkinter import filedialog, messagebox, scrolledtext
	import customtkinter as ctk
	
	# Класс правила Маркова
	class MarkovRules:
	def __init__(self, left, right, stop=False):
	self.left = left # левая чсть правила
	self.right = right # правая часть правила
	self.stop = stop # правило с точкой или нет bool типа
	
	# поиск левой части в строке и замена её на правую(только первое вхождение)
	def apply(self, text):
	if self.left in text:
	new_text = text.replace(self.left, self.right, 1)
	return new_text, self.stop
	return None, False
	
	
	# Класс системы Маркова
	class MarkovSystem:
	def __init__(self, rules, animation=False):
	self.rules = rules # Список правил (Объекты касса MarkovRules)
	self.steps = 0 # счётчик замен
	self.applied_rules = {} # для статистики: словарь хранит в себе количество применений каждого правила
	self.lengths = [] # хранит длину строки на каждом шагу
	self.animation = animation # bool - нужна ли анимация или нет
	
	# запуск системы
	def run(self, text):
	print("\n🔹 Запуск системы Маркова...\n")
	log = [] # список всех шагов преобразований
	
	while True:
	#проверка правил по порядку
	for rule in self.rules:
	new_text, stop = rule.apply(text)
	if new_text: # если не None(есть совпадение)
	self.steps += 1
	
	# Подсчёт количества применений каждого правила
	key = f"{rule.left} → {rule.right}"
	self.applied_rules[key] = self.applied_rules.get(key, 0) + 1 
	
	self.lengths.append(len(new_text)) # длина новой строки
	visual = f"{text} → {new_text}" # текстовое представление шага
	print(f"({self.steps}) {visual}  | правило: {key}") # текстовое представление шага в консоли
	log.append(visual) # запоминаем в логах
	text = new_text # обновляем строку
	
	# имитация анимации
	if self.animation:
	time.sleep(0.6)
	
	# если stop не False
	if stop:
	print("\n Алгоритм завершён (правило с точкой).")
	self.save_log(log)
	self.show_stats()
	return text
	
	# После любого успешного шага — снова начинаем с первого правила
	break
	else:
	# Если не нашлось ни одного применимого правила
	print("\n Больше нет применимых правил.")
	self.save_log(log)
	self.show_stats()
	return text
	
	# функция сохранения журнала шагов в log.txt
	@staticmethod
	def save_log(log):
	with open("C:\\Users\\user\\desktop\\log.txt", "w", encoding="utf-8") as f:
	f.write("\n".join(log))
	print(" Журнал шагов сохранён в log.txt\n")
	
	# Показ статистики
	def show_stats(self):
	print(" Статистика выполнения:\n")
	print(f"• Всего шагов: {self.steps}")
	
	if self.lengths:
	# рассчитываем среднюю длину строки
	avg_len = sum(self.lengths) / len(self.lengths)
	print(f"• Средняя длина строки: {avg_len:.2f}")
	
	if self.applied_rules:
	# находим правило, которое применяеся чаще всего
	most_common_rule = max(self.applied_rules, key=self.applied_rules.get)
	count = self.applied_rules[most_common_rule]
	print(f"• Самое частое правило: {most_common_rule} ({count} раз)")
	
	
	
	# Функция загрузки правил из файла
	def load_rules_from_file():
	base_path = "C:\\Users\\user\\Desktop\\"
	filename = input("Введите имя файла с правилами (например: rules.txt): ").strip()
	full_path = os.path.join(base_path, filename)
	
	rules = [] # павила
	
	# проверка существования файла
	if not os.path.exists(full_path):
	print(f"⚠️ Файл '{full_path}' не найден.")
	return rules
	# читаем файл и извлекаем правила
	with open(full_path, "r", encoding="utf-8") as f:
	for line in f:
	line = line.strip()
	if not line or "->" not in line:
	continue
	left, right = line.split("->")
	left, right = left.strip(), right.strip()
	stop = right.endswith(".")
	if stop:
	right = right[:-1]
	
	# создаём объект правила и добавляет в список
	rules.append(MarkovRules(left, right, stop))
	
	print(f"✅ Правила успешно загружены из файла: {filename}\n")
	return rules
	
	# Функция ручного ввода правил
	def input_rules():
	print("\nВвод правил вручную (используйте формат 'левая_часть -> правая_часть')")
	print("Чтобы завершить ввод, просто нажмите Enter на пустой строке.")
	
	rules = []  # Список правил, который будем заполнять
	
	
	while True: 
	line = input("Правило (пример: ab -> c или a -> x.): ").strip()
	
	# Если пустая строка, выходим из цикла
	if not line:
	break
	
	# проверка формата правила
	if "->" not in line: 
	print("Ошибка: правило должно содержать '->'")
	continue
	
	left, right = line.split("->")
	left, right = left.strip(), right.strip()  # Убираем лишние пробелы
	
	stop = right.endswith(".")  # если заканчивается точкой правила завершающее
	if stop:
	right = right[:-1]  # убираем точку из правой части
	
	rules.append(MarkovRules(left, right, stop))  # создаём объект правила и добавляем в список
	
	return rules
	
	
	def start_gui():
	ctk.set_appearance_mode("dark")  # "light", "dark", "system"
	ctk.set_default_color_theme("blue")
	
	app = ctk.CTk()
	app.title("Система Маркова")
	app.geometry("900x650")
	
	# Заголовок
	ctk.CTkLabel(app, text="💡 Система Маркова", font=("Segoe UI", 24, "bold")).pack(pady=15)
	
	# Основной контейнер
	frame = ctk.CTkFrame(app, corner_radius=12)
	frame.pack(pady=10, padx=20, fill="both", expand=True)
	
	# Исходная строка
	input_entry = ctk.CTkEntry(frame, placeholder_text="Введите исходную строку...", height=40)
	input_entry.pack(pady=15, padx=20, fill="x")
	
	# Кнопки
	button_frame = ctk.CTkFrame(frame, fg_color="transparent")
	button_frame.pack(pady=5)
	
	rules_data = []
	anim_var = tk.BooleanVar(value=False)
	
	# 📂 Загрузка правил из файла
	def load_rules():
	nonlocal rules_data
	path = filedialog.askopenfilename(title="Выберите файл с правилами",
	filetypes=[("Текстовые файлы", "*.txt")])
	if not path:
	return
	
	try:
	rules_data.clear()
	with open(path, "r", encoding="utf-8") as f:
	lines = f.readlines()
	
	output_box.delete("1.0", "end")
	output_box.insert("end", f"📘 Загружен файл: {path}\n\n")
	for line in lines:
	line = line.strip()
	if not line or "->" not in line:
	continue
	left, right = line.split("->")
	left, right = left.strip(), right.strip()
	stop = right.endswith(".")
	if stop:
	right = right[:-1]
	rules_data.append(MarkovRules(left, right, stop))
	output_box.insert("end", f"{left} → {right}{'.' if stop else ''}\n")
	output_box.insert("end", "\n✅ Правила успешно загружены.\n")
	except Exception as e:
	messagebox.showerror("Ошибка", f"Не удалось загрузить файл:\n{e}")
	
	#Ручной ввод правил
	def input_rules_gui():
	def save_rules():
	rules_data.clear()
	lines = text_area.get("1.0", "end").strip().splitlines()
	for line in lines:
	if "->" in line:
	left, right = line.split("->")
	left, right = left.strip(), right.strip()
	stop = right.endswith(".")
	if stop:
	right = right[:-1]
	rules_data.append(MarkovRules(left, right, stop))
	top.destroy()
	output_box.insert("end", f"\n📝 Добавлено {len(rules_data)} правил из ручного ввода.\n")
	
	top = ctk.CTkToplevel(app)
	top.title("Ручной ввод правил")
	
	# Делает окно привязанным к главному и всегда поверх
	top.transient(app)
	top.grab_set()
	top.focus_set()
	
	tk.Label(top, text="Введите правила (формат: a -> b или x -> y.)").pack(pady=5)
	text_area = scrolledtext.ScrolledText(top, width=50, height=15)
	text_area.pack(padx=10, pady=5)
	ctk.CTkButton(top, text="Сохранить", command=save_rules).pack(pady=5)
	
	def run_markov():
	if not rules_data:
	messagebox.showwarning("Нет правил", "Сначала загрузите или введите правила!")
	return
	text = input_entry.get().strip()
	if not text:
	messagebox.showwarning("Пустая строка", "Введите исходный текст.")
	return
	
	output_box.delete("1.0", "end")
	system = MarkovSystem(rules_data, animation=False)
	log = []
	
	def show_stats_and_result(final_text):
	system.save_log(log)
	stats = f"\n📊 Статистика:\n• Всего шагов: {system.steps}\n"
	if system.lengths:
	stats += f"• Средняя длина строки: {sum(system.lengths)/len(system.lengths):.2f}\n"
	if system.applied_rules:
	most_common_rule = max(system.applied_rules, key=system.applied_rules.get)
	count = system.applied_rules[most_common_rule]
	stats += f"• Самое частое правило: {most_common_rule} ({count} раз)\n"
	stats += f"\n🏁 Результат: {final_text}\n"
	output_box.insert("end", stats)
	output_box.see("end")
	
	if anim_var.get():  # Анимация включена
	def step_animation(current_text):
	for rule in system.rules:
	new_text, stop = rule.apply(current_text)
	if new_text:
	system.steps += 1
	key = f"{rule.left} → {rule.right}"
	system.applied_rules[key] = system.applied_rules.get(key, 0) + 1
	system.lengths.append(len(new_text))
	
	visual = f"({system.steps}) {current_text} → {new_text}  | правило: {key}"
	log.append(visual)
	output_box.insert("end", visual + "\n")
	output_box.see("end")
	
	if stop:
	show_stats_and_result(new_text)
	return
	
	# Переходим к следующему шагу через after
	app.after(300, lambda c=new_text: step_animation(c))
	return
	else:
	# Нет больше применимых правил
	show_stats_and_result(current_text)
	
	step_animation(text)
	
	else:  # Анимация выключена — все сразу
	current_text = text
	while True:
	for rule in system.rules:
	new_text, stop = rule.apply(current_text)
	if new_text:
	system.steps += 1
	key = f"{rule.left} → {rule.right}"
	system.applied_rules[key] = system.applied_rules.get(key, 0) + 1
	system.lengths.append(len(new_text))
	
	visual = f"({system.steps}) {current_text} → {new_text}  | правило: {key}"
	log.append(visual)
	output_box.insert("end", visual + "\n")
	output_box.see("end")
	
	current_text = new_text
	if stop:
	show_stats_and_result(current_text)
	return
	break
	else:
	show_stats_and_result(current_text)
	return
	
	# Кнопки
	ctk.CTkButton(button_frame, text="📂 Загрузить правила", command=load_rules, width=180).grid(row=0, column=0, padx=5)
	ctk.CTkButton(button_frame, text="📝 Ручной ввод", command=input_rules_gui, width=180).grid(row=0, column=1, padx=5)
	ctk.CTkCheckBox(button_frame, text="Анимация", variable=anim_var).grid(row=0, column=2, padx=5)
	ctk.CTkButton(button_frame, text="▶️ Запустить", command=run_markov, width=180).grid(row=0, column=3, padx=5)
	
	# Поле вывода
	output_box = ctk.CTkTextbox(frame, height=350)
	output_box.pack(pady=10, padx=20, fill="both", expand=True)
	output_box.insert("end", "Шаги выполнения алгоритма...\n")
	
	app.mainloop()
	
	if __name__ == "__main__":
	start_gui()
	# # Основной цикл
	# def main():
	#     while True:
	#         choice = main_menu() # шоу меню
	#         if choice == "1": # выбор загрузки правил из файла
	#             rules = load_rules_from_file() 
	#         elif choice == "2": # пользователь выбрал ручной ввод правил
	#             rules = input_rules()
	#         elif choice == "3": # выход
	#             print("👋 Завершение работы.")
	#             break
	#         else: # некорректный ввод
	#             print("❌ Неверный выбор.\n")
	#             continue
	
	#         # проверка: есть ли правила
	#         if not rules:
	#             print("⚠️ Нет загруженных правил.")
	#             continue
	
	#         text = input("\nВведите исходную строку: ").strip()
	
	#         anim_choice = input("Включить анимацию (y/n)? ").lower() == "y"
	
	#         #запуск
	#         system = MarkovSystem(rules, animation=anim_choice)
	#         result = system.run(text)
	
	#         print(f"\n🏁 Результат: {result}\n")
	#         input("Нажмите Enter, чтобы вернуться в меню...")
	
	
	# if __name__ == "__main__":
	#     main()
	
\end{lstlisting}